\section{Observation Knowledge Is Learned but Not Used}
\label{sec:mechanism}

\begin{figure}[t]
\centering
\figplaceholder{figures/fig3_mechanism.pdf}{5.2cm}
\caption{\textbf{(a)}~Held-out action KL of each SFT arm to \armA{}. The grey bar and dotted line mark the data-order noise floor: an action-only model trained on the same data in a different order. \textbf{(b)}~Observation cross-entropy (nats) by part, on unseen (held-out) databases and on training databases. \armB{} learns the template and the result header everywhere, but predicts row values better only on databases it was trained on.}
\label{fig:mechanism}
\end{figure}

\subsection{The action policy is unchanged}
\label{sec:equiv}
We measure the mean per-token KL between action distributions on 300 held-out trajectories from databases that SFT never saw. For scale, SFT moves the policy far from the base model: the KL from Qwen3-4B is 0.364 for \armA{} and 0.363 for \armB{}. Against that, \armB{} stays very close to \armA{}, with $\KL(\armB\|\armA)=0.0049$ (Figure~\ref{fig:mechanism}a).

This is \emph{less} than the KL between two action-only models trained on the same data in a different order (0.0070). It is of the same order as the effect of a 15\% lower learning rate (\armC{}, 0.0010) or of Wikipedia dilution (\armD{}, 0.0027).

Behaviour after observations agrees. On BIRD before RL, \armA{} and \armB{} produce a similar share of trajectories with an SQL error (16.5\% vs.\ 15.7\%). They also solve the same share of those trajectories (30.3\% in both). After GRPO the figures are 13.5\% vs.\ 13.6\%, and 37.7\% vs.\ 37.5\%.

GRPO does change how the agent uses observations. Re-issuing a query that has just failed drops from 23--26\% of failures to 6--7\%. But it changes \armA{} and \armB{} by the same amount (App.~\ref{app:behaviour}).

\subsection{What the model learns}
\label{sec:learned}
\armB{} does learn to predict observations. Its observation cross-entropy on training trajectories falls from 3.57 to 0.25 within the epoch. Most of this learning, however, does not transfer.

We split observation tokens into three parts: the fixed template (\texttt{<observation>} tags and the turn reminder), the result header (the first line of a result table, which encodes its column count), and the row values. We then compare databases that SFT never saw with databases it trained on (Figure~\ref{fig:mechanism}b; full numbers in App.~\ref{app:obsce}). Structure transfers: on unseen databases \armB{} predicts the template almost perfectly (CE 0.03 vs.\ 2.79 for \armA{}), and its CE on the result header falls from 5.64 to 1.02 on the first tokens and from 2.22 to 0.30 on the rest of the line. Content does not. On unseen databases, row values are predicted no better than by \armA{} (1.92 vs.\ 1.88 over 35,845 tokens), whereas on training databases \armB{} fits them closely (0.70 vs.\ 1.60). Error messages and empty results are too rare in held-out teacher trajectories (5 and 3 segments) to analyse.

\armB{} thus gains two kinds of knowledge. One is structure it can predict but has no reason to act on, since the template and header follow from the query it has just written. The other is database content that cannot transfer to the unseen databases it is evaluated on. Neither tells the agent whether its current result answers the question.

The gradients agree. Action and observation gradients are nearly orthogonal at the end of SFT (cosine 0.007), as \citet{zhang2026actobs} also report. On training trajectories, the observation gradient is no larger than the action gradient: the norm ratio is 1.09, against 5.49 on held-out trajectories.

The pre-registered dilution account also fails at the SFT stage: \armB{} moves the action policy as much as \armA{} does, although its pooled loss gives actions only 69\% of each update. App.~\ref{app:share} discusses why.
